% Adición focal propuesta; no sustituye el álgebra electrónica existente.
% Inserción: después de la proposición «Álgebra del bloque electrónico local»
% y su párrafo explicativo, antes de «Transferencia entre hojas...».
% Procedencia y alcance: LOCALIZADORES_Y_ALCANCE.md, en esta carpeta.
\subsection{Représentation spinorielle du bloc électronique et hélicité}
\label{sec:el-representacion-espinorial}

La structure électronique possède une représentation qui conserve davantage
d'information que son évaluation scalaire. Le plan principal des deux
projecteurs d'APP tridimensionnelle, sélectionné par le mode tritique
\((1,-1,0)^3\), a déjà fourni les opérateurs \(S\) et \(J\).
Le premier est une involution symétrique ; le second est le commutateur
normalisé et conserve l'orientation des feuillets. Leur composition permet
de construire explicitement une représentation des rotations et de séparer ses
deux composantes d'hélicité. Cette construction agit sur le bloc
électronique ; le registre TPK de la route et sa mémoire demeurent les données
de provenance du bloc.

Soit \(E_e\) le plan réel principal, muni du produit scalaire hérité
des espérances conditionnelles, et soit
\(\mathscr S_e=E_e\otimes_{\mathbb R}\mathbb C\) sa complexification.
Les opérateurs déjà obtenus satisfont
\[
 S^*=S,\qquad J^*=-J,\qquad
 S^2=I,\qquad J^2=-I,\qquad SJ=-JS.
\]
La complexification étend le corps de représentation. Elle conserve les
opérateurs réels et permet de distinguer l'unité scalaire \(i\) de
l'endomorphisme \(J\), bien que tous deux aient pour carré \(-1\) dans leurs
domaines respectifs.

\begin{proposition}[Triplet hermitien induit par les deux feuillets]
\label{prop:el-terna-hermitica}
Les opérateurs
\begin{equation}
 \sigma_1=SJ,\qquad \sigma_2=-iJ,\qquad \sigma_3=S
 \label{eq:el-terna-hermitica}
\end{equation}
sont hermitiens et de trace nulle. Ils satisfont
\begin{equation}
 \sigma_a\sigma_b
 =\delta_{ab}I+i\sum_{c=1}^3\epsilon_{abc}\sigma_c,
 \qquad
 \frac12\operatorname{tr}(\sigma_a\sigma_b)=\delta_{ab},
 \label{eq:el-algebra-terna}
\end{equation}
où \(\epsilon_{123}=1\). Ils constituent donc une base orthonormée
de l'espace réel \(V_e\) des endomorphismes hermitiens de trace nulle de
\(\mathscr S_e\), pour le produit
\(\langle X,Y\rangle=\frac12\operatorname{tr}(XY)\).
\end{proposition}
\begin{proof}
Les relations précédentes donnent \((SJ)^*=(-J)S=SJ\) et
\((-iJ)^*=-iJ\). Les carrés des trois matrices valent \(I\).
De plus,
\[
 (SJ)(-iJ)=iS,\qquad (-iJ)S=iSJ,\qquad S(SJ)=J=i(-iJ).
\]
Inverser l'ordre change le signe. Cela établit la première identité
de \eqref{eq:el-algebra-terna}. Dans la base orthonormée du plan principal,
les matrices sont
\[
 \sigma_1=\begin{pmatrix}0&1\\1&0\end{pmatrix},\quad
 \sigma_2=\begin{pmatrix}0&-i\\i&0\end{pmatrix},\quad
 \sigma_3=\begin{pmatrix}1&0\\0&-1\end{pmatrix}.
\]
Leur trace est nulle, et prendre les traces dans l'identité de produit établit
l'orthonormalité. Tout endomorphisme hermitien de trace nulle a la forme
\(\bigl(\begin{smallmatrix}z&x-iy\\x+iy&-z\end{smallmatrix}\bigr)\),
avec \(x,y,z\) réels ; il appartient donc à leur espace vectoriel réel engendré.
\end{proof}

Dans la représentation matricielle usuelle, ce triplet porte le nom de
matrices de Pauli. Son utilisation ici est postérieure à la construction de
\(P,Q,S,J\) : il s'obtient par les produits
\eqref{eq:el-terna-hermitica}. L'espace des directions de cette
représentation est la sphère unité de \(V_e\), dont la forme positive
vient d'être construite au moyen de la trace.

\begin{proposition}[Générateurs de rotation et décomposition directionnelle]
\label{prop:el-helicidad}
Définissons \(L_a=\sigma_a/2\). Alors
\begin{equation}
 [L_a,L_b]=i\sum_c\epsilon_{abc}L_c,
 \qquad \sum_{a=1}^3L_a^2=\frac34 I.
 \label{eq:el-generadores-spin}
\end{equation}
Pour \(n=(n_1,n_2,n_3)\) avec \(\sum_an_a^2=1\), posons
\[
 H_e(n)=\sum_an_a\sigma_a,\qquad
 h_e(n)=\frac12H_e(n),\qquad
 \Pi_\pm(n)=\frac12\bigl(I\pm H_e(n)\bigr).
\]
La formule linéaire \(H_e(v)=\sum_av_a\sigma_a\) est également utilisée
pour \(v\) de norme arbitraire.
On a
\begin{equation}
 \mathscr S_e=\operatorname{im}\Pi_+(n)
             \mathbin{\oplus}\operatorname{im}\Pi_-(n),
 \qquad
 h_e(n)\Pi_\pm(n)=\pm\frac12\Pi_\pm(n).
 \label{eq:el-descomposicion-helicidad}
\end{equation}
Les deux composantes sont des droites complexes orthogonales. Le transport de
rotation autour de \(n\) est
\begin{equation}
 U_n(\theta)=\exp\!\left(-\frac{i\theta}{2}H_e(n)\right)
 =\cos\frac\theta2\,I-i\sin\frac\theta2\,H_e(n),
 \qquad
 U_n(2\pi)=-I,\quad U_n(4\pi)=I.
 \label{eq:el-retorno-spin}
\end{equation}
\end{proposition}
\begin{proof}
L'identité de produit du triplet donne les commutateurs de
\eqref{eq:el-generadores-spin} ; la somme de leurs carrés vaut \(3I/4\).
Les termes antisymétriques s'annulent lorsqu'on multiplie
\(H_e(n)\) par lui-même, d'où \(H_e(n)^2=I\).
On en déduit
\[
 \Pi_\pm(n)^2=\Pi_\pm(n),\qquad
 \Pi_+(n)\Pi_-(n)=0,\qquad \Pi_+(n)+\Pi_-(n)=I.
\]
Ils sont hermitiens et de trace \(1\) ; leur rang vaut donc \(1\).
Multiplier par \(H_e(n)/2\) établit
\eqref{eq:el-descomposicion-helicidad}. La série exponentielle se
sépare en puissances paires et impaires grâce à \(H_e(n)^2=I\), ce qui
produit \eqref{eq:el-retorno-spin}. Elle établit également directement
l'unitarité et la loi \(U_n(\theta+\psi)=U_n(\theta)U_n(\psi)\).

La normalisation angulaire se vérifie sur \(V_e\). Pour tout
\(v\in\mathbb R^3\), la même identité de produit donne
\[
 U_n(\theta)H_e(v)U_n(\theta)^*
 =H_e\!\left(v\cos\theta+(n\mathbin{\times}v)\sin\theta
       +n(n\mathbin{\cdot}v)(1-\cos\theta)\right).
\]
L'action induite est une rotation d'angle \(\theta\) et d'axe \(n\).
Le facteur \(1/2\) dans le générateur spinoriel est donc celui qui
correspond à cette paramétrisation angulaire. Une rotation complète rétablit
la direction dans \(V_e\) et change le signe du vecteur de
\(\mathscr S_e\) ; deux rotations complètes rétablissent les deux.
\end{proof}

La représentation est irréductible : un sous-espace complexe invariant sous
les trois matrices serait invariant sous toute l'algèbre \(M_2(\mathbb C)\),
qu'elles engendrent. Ses seuls sous-espaces invariants sont \(0\) et
\(\mathscr S_e\). Les relations
\eqref{eq:el-generadores-spin} identifient ainsi la représentation de
spin \(1/2\). Si \(n\) est la direction de propagation dans une
réalisation cinématique, \(h_e(n)\) est son hélicité sans dimension, de
valeurs \(\pm1/2\). La direction cinématique est déclarée lors de cette
réalisation ; les deux projecteurs sont définis pour chaque direction
avant la sélection d'un état de propagation particulier.

\paragraph{Conservation de la structure électronique et types d'inversion.}
L'inversion de direction possède une loi exacte :
\[
 H_e(-n)=-H_e(n),\qquad \Pi_\pm(-n)=\Pi_\mp(n).
\]
Cette opération échange les étiquettes d'hélicité relatives à une
direction. La conjugaison de charge agit, en revanche, sur la signature
orientée de la route matérielle ; l'inversion cellulaire agit sur les
coordonnées de l'atlas. Leurs domaines et leurs actions demeurent distincts.
L'unicité de la cellule \((5,5)\) est compatible avec les deux droites de
\eqref{eq:el-descomposicion-helicidad} : une cellule de l'atlas n'est pas un
vecteur de la représentation spinorielle qu'elle porte.

L'évaluation électronique \(\varepsilon_e^\beta\) conserve sa
formule et ses opérateurs de lecture. Sur \(\mathscr S_e\), sa réalisation scalaire
est \(\varepsilon_e^\beta I\), qui commute avec les générateurs et avec
\(\Pi_\pm(n)\). La décomposition d'hélicité ajoute l'information
rotationnelle sans modifier l'évaluation ni convertir le spin en un
coefficient multiplicatif de la masse.

La chiralité correspond à la graduation de la représentation relativiste
utilisée : une involution \(\Gamma_{\rm ch}\) détermine ses
projecteurs \((I\mp\Gamma_{\rm ch})/2\). L'hélicité de
\eqref{eq:el-descomposicion-helicidad} dépend en outre de \(n\).
L'holonomie d'un fibré réel de Möbius décrit, quant à elle, le changement
de signe d'une droite lorsqu'on parcourt un lacet de sa base. La loi de cette droite,
le retour spinoriel \eqref{eq:el-retorno-spin} et l'holonomie nonadique
de phase avec avance de mémoire sont des lois de transport sur des objets
spécifiés. Leur comparaison conserve ces objets : la forme hélicoïdale
d'une trajectoire ne remplace pas la représentation qui vient d'être
construite, ni ne fixe à elle seule la chiralité ou la charge.
